% !TEX program = xelatex
% =============================================================================
% cn-letter · 求职自荐信（高管版，应聘 CEO）
%
% 样张为应聘首席执行官岗位的高管自荐信，致董事会及提名委员会。结构：
% 应聘岗位与渠道 → 三段匹配点（经营业绩/组织与治理/行业判断——每段
% 一个主题，用可核验的数字与事实支撑，不堆形容词）→ 面谈期待与附件
% 说明 → 落款附联系方式。普通岗位自荐信同构，换称谓与事实即可，
% 与 templates/resume 简历配套投递。
%
% 编译：bash scripts/compile.sh templates/cn-letter/cover-letter.tex --preview
% =============================================================================
\documentclass[zihao=-4]{ctexart}
\usepackage{letter}

\begin{document}

\ltTitle{自\hspace{1\ccwd}荐\hspace{1\ccwd}信}

\ltDear{尊敬的董事会及提名委员会：}

获悉贵公司面向社会公开遴选\hl{首席执行官（CEO）}，本人经对遴选
条件逐项对照后，特此自荐。

经营业绩方面：本人现任示例科技股份有限公司\hl{首席运营官（COO）}，
全面负责公司经营管理。任职四年间，公司营业收入自3.2亿元增长至12.6亿
元，年复合增长率约40\%；主导完成B轮、C轮合计15亿元股权融资，并推动
公司进入首次公开发行股票的申报程序。

组织与治理方面：本人主持搭建了自200人至800人规模的组织体系与
职业经理人梯队，建立全面预算、内控合规与信息披露制度；长期直接向
董事会汇报，具备与投资人、监管机构及中介机构沟通的成熟经验。

行业判断方面：本人深耕人工智能行业十二年，主导过两次关键技术
路线切换，对行业周期与商业化路径有经过验证的判断框架，与贵公司
新五年战略中的智能化转型方向高度契合。

本人期待有机会向董事会当面汇报经营思路与就任后的百日计划。随信
附上个人履历及可供核验的业绩证明材料清单，敬请审阅。

感谢董事会及提名委员会拨冗垂阅！

\ltClosing

\ltSign{自荐人：林××}
\ltSign{联系电话：138\,0000\,0000\quad 电子邮箱：linxx@example.com}
\ltSign{2026年7月18日}

\end{document}
